\chapter{Future Work: ML-Based Anomaly Detection}\label{ch:future-ml}

This chapter documents the roadmap for transitioning from the current
\textbf{statistical anomaly detection} methods (3$\sigma$ thresholds, MAD,
percentile-based rules) to \textbf{machine learning-based} approaches.

\textit{Note: The current implementation uses statistical process control
techniques. The term ``ML-based'' has been removed from the main specification
to accurately reflect the current state. This chapter describes v2.0 plans.}

\section{Current state: Statistical detection}

The v1.0 implementation uses four statistical checks (see Chapter~3):
\begin{enumerate}
  \item Period spike detection ($>3\sigma$ from 12-month rolling mean)
  \item Sign flip detection (balance sign reversal)
  \item New account detection (no prior history)
  \item Stale balance detection (unchanged $>$6 months)
\end{enumerate}

These methods are effective for \textbf{known anomaly patterns} but have
limitations:
\begin{itemize}
  \item Cannot learn new anomaly patterns from historical data
  \item Require manual threshold tuning per entity
  \item Assume statistical stationarity (violated during restructuring)
  \item Cannot capture multi-variate correlations
\end{itemize}

\section{v2.0 roadmap: Machine learning approaches}

\subsection{Isolation Forest}

\begin{table}[htbp]
\centering
\caption{Isolation Forest specification.}
\footnotesize
\begin{tabularx}{\linewidth}{@{}l X@{}}
\toprule
\textbf{Attribute} & \textbf{Value} \\ \midrule
Algorithm & Isolation Forest (scikit-learn) \\
Input features & 12-month balance history, variance ratios, sign changes \\
Training data & 2 years historical TB data per entity \\
Contamination parameter & 0.05 (5\% expected anomaly rate) \\
Number of estimators & 100 trees \\
Retraining frequency & Quarterly \\
\bottomrule
\end{tabularx}
\end{table}

\subsection{Autoencoder}

\begin{table}[htbp]
\centering
\caption{Autoencoder specification.}
\footnotesize
\begin{tabularx}{\linewidth}{@{}l X@{}}
\toprule
\textbf{Attribute} & \textbf{Value} \\ \midrule
Architecture & Dense autoencoder (128-64-32-64-128) \\
Input features & Normalized balance vectors (12 months) \\
Training data & 2 years historical TB data, anomalies removed \\
Loss function & Mean squared error \\
Anomaly threshold & Reconstruction error $>$ 95th percentile \\
Retraining frequency & Monthly \\
\bottomrule
\end{tabularx}
\end{table}

\subsection{Time-series forecasting}

For accounts with seasonal patterns:
\begin{itemize}
  \item Prophet model for trend + seasonality decomposition
  \item Anomaly = actual deviates from forecast $>$2$\sigma$
  \item Requires $\geq$24 months history
\end{itemize}

\section{Transition plan}

\subsection{Phase 1: Shadow mode (Months 1--3)}

\begin{enumerate}
  \item Deploy ML models alongside statistical checks
  \item Log ML predictions without surfacing to users
  \item Compare ML vs. statistical detection rates
  \item Measure precision/recall on human-labelled anomalies
\end{enumerate}

\subsection{Phase 2: Hybrid mode (Months 4--6)}

\begin{enumerate}
  \item Use statistical checks as primary
  \item Add ML-detected anomalies as ``additional flags''
  \item Collect user feedback on ML-only detections
  \item Tune thresholds based on feedback
\end{enumerate}

\subsection{Phase 3: ML primary (Month 7+)}

\begin{enumerate}
  \item Switch to ML as primary detection
  \item Keep statistical checks as fallback
  \item Establish continuous monitoring of model drift
  \item Quarterly model retraining and validation
\end{enumerate}

\section{Success criteria}

The transition is complete when:
\begin{itemize}
  \item ML precision $\geq$ 85\% (verified anomalies / flagged anomalies)
  \item ML recall $\geq$ 90\% (detected anomalies / actual anomalies)
  \item False positive rate $\leq$ 10\%
  \item Human review load reduced by $\geq$ 20\%
  \item Model retraining automated and validated
\end{itemize}

\section{Data requirements}

\begin{table}[htbp]
\centering
\caption{Training data requirements for ML models.}
\footnotesize
\begin{tabularx}{\linewidth}{@{}l r X@{}}
\toprule
\textbf{Data Type} & \textbf{Volume} & \textbf{Source} \\ \midrule
Historical TB extracts & 24 months & SAP archives \\
Labelled anomalies & 500+ per entity type & Human review logs \\
Seasonal patterns & 36 months & For Prophet model \\
Entity metadata & All 234 entities & Master data \\
\bottomrule
\end{tabularx}
\end{table}

\section{Timeline}

\begin{table}[htbp]
\centering
\caption{ML transition timeline (conditional, not committed).}
\label{tab:ml-timeline}
\footnotesize
\begin{tabularx}{\linewidth}{@{}l l X@{}}
\toprule
\textbf{Milestone} & \textbf{Target window} & \textbf{Gating dependencies} \\ \midrule
Data collection complete & Q3--Q4 2026 & Historical data extraction; labelled anomaly backlog $\ge$500/entity type \\
Isolation Forest prototype & Q1 2027 & Data collection \emph{and} 2+ quarters of dashboard audit trail (for labels) \\
Shadow mode deployment & Q2 2027 & Prototype passes offline validation on held-out set \\
Hybrid mode pilot (HKG only) & Q3 2027 & Shadow mode shows $\ge$90\% agreement with statistical checks on true positives \\
ML primary (limited rollout) & Q4 2027--Q1 2028 & Hybrid success on HKG + 2 additional LEs \\
Full production & Q2--Q4 2028 & All onboarded entities validated (paced by LE onboarding in~\cref{sec:params-234}) \\
\bottomrule
\end{tabularx}
\end{table}

\paragraph{Schedule realism.}
The dates in Table~\ref{tab:ml-timeline} are \emph{target windows},
not committed dates. The v1.0 implementation has no ML model today
and no labelled anomaly corpus; the earliest honest shadow-mode
deployment is therefore gated on accumulating $\ge$2~quarters of
production dashboard decisions (start of the corpus coincides with
v1.0 go-live). Each milestone's entry is a \textbf{target}; missing
a target by one quarter does not block the v1.0 programme and does
not downgrade the v1.0 statistical methods. The programme
management office reviews this table quarterly and revises windows
based on: (a) labelled corpus volume, (b) LE onboarding pace
(\cref{sec:params-234}), and (c) controller feedback on
statistical detection false-positive rate. A formal decision to
commit to shadow-mode deployment is taken by the TB programme
board one quarter before the target.